% \iffalse
%<*driver>
\ProvidesFile{seqenseigne-theme.dtx}%
[2025/10/24 v1.0-dev seqenseigne – Module theme (couleurs, styles, visuels)]
%</driver>
% \fi
%
% =============================================================
% seqenseigne-theme.dtx
% Module : THEME
% Rôle   :
%   - palette de couleurs (toutes les \definecolor et macros
%     \seqCouleurXXX)
%   - réglages graphiques réutilisables :
%       cadres colorés, icônes info/objectif/cours, encadrés objectifs
%   - styles des tableaux (largeurs de colonnes L{...}, etc.)
%   - mise en forme visuelle des environnements pédagogiques
%     (ex : fond coloré, bandeau, pictos, etc.)
%
%   → Tout ce qui, dans ton dtx actuel, gère l'apparence
%     plus que la logique.
%
% =============================================================
%
%<*package>
\NeedsTeXFormat{LaTeX2e}%
\ProvidesPackage{seqenseigne-theme}%
[2025/10/24 v1.0-dev seqenseigne theme module]

% -------------------------------------------------------------
% Dépendances purement graphiques
% -------------------------------------------------------------
\RequirePackage{adjustbox}%
\RequirePackage[theorems,breakable,skins]{tcolorbox}%
\RequirePackage{xcolor}%
\RequirePackage{tikz}%
\RequirePackage{adjustbox}%
\RequirePackage{multirow}%
\RequirePackage{colortbl}%
\RequirePackage{datetime}%
\RequirePackage{geometry}%
\RequirePackage{fancyhdr}%
\RequirePackage{multicol}%
\RequirePackage{pifont}%
\RequirePackage[inline]{enumitem}%
\RequirePackage{xkeyval}%
\RequirePackage{xifthen}%

%% ------------------------------------------------------------------
%% ---- Traitement des options --------------------------------------
%% ------------------------------------------------------------------
\newif\ifdocstd\docstdtrue
\DeclareOption{ltxdoc}{\docstdfalse}%
\ProcessOptions\relax

% -------------------------------------------------------------
% Palette et variables graphiques
% -------------------------------------------------------------
% \begin{macro}{\seqCouleurPrimaire}%
% \IndexMacro{\seqCouleurPrimaire}%
% etc.
%
% \begin{macrocode}%
% \definecolor{seqCouleurPrimaire}{RGB}{...}%
% \newcommand{\seqCouleurTexte}{seqCouleurPrimaire}%
% \colorlet{seqCouleurTitre}{...}%
% \newcommand{\seqCouleurObjectif}{...}%
% \newcommand{\seqCouleurAvert}{...}%
% \newcommand{\seqCouleurInfo}{...}%
% \newcommand{\seqCouleurFondBloc}{...}%
% \newcommand{\seqCouleurNiveauF}{...} % etc pour F / A / E
% \newcommand{\seqCouleurNiveauA}{...}%
% \newcommand{\seqCouleurNiveauE}{...}%
% \end{macrocode}%
%
%% ------------------------------------------------------------------
%% ---- Définition des couleurs -------------------------------------
%% ------------------------------------------------------------------
%% ---- Nombres et calcul
\newcommand{\seqCouleurRougeVif}{red!85!black}%
\newcommand{\seqCouleurRougeModere}{red!50!white}%
\newcommand{\seqCouleurRougeMoyen}{red!25!white}%
\newcommand{\seqCouleurRougePale}{red!10!white}%
\newcommand{\seqCouleurJauneModere}{yellow!50!white}%
\newcommand{\seqCouleurJaunePale}{yellow!10!white}%
%% ---- Organisation des données et fonction
\newcommand{\seqCouleurBleuVif}{blue!85!black}%
\newcommand{\seqCouleurBleuModere}{blue!50!white}%
\newcommand{\seqCouleurBleuMoyen}{blue!25!white}%
\newcommand{\seqCouleurBleuPale}{blue!10!white}%
\newcommand{\seqCouleurCanardModere}{teal!50!white}%
\newcommand{\seqCouleurCanardPale}{teal!10!white}%
%% ---- Grandeurs et mesures
\newcommand{\seqCouleurMarronVif}{brown!85!black}%
\newcommand{\seqCouleurMarronModere}{brown!50!white}%
\newcommand{\seqCouleurMarronMoyen}{brown!25!white}%
\newcommand{\seqCouleurMarronPale}{brown!10!white}%
\newcommand{\seqCouleurOrangeModere}{orange!50!white}%
\newcommand{\seqCouleurOrangePale}{orange!10!white}%
%% ---- Espace et géométrie
\newcommand{\seqCouleurVioletVif}{violet!85!black}%
\newcommand{\seqCouleurVioletMoyen}{violet!25!white}%
\newcommand{\seqCouleurVioletModere}{violet!50!white}%
\newcommand{\seqCouleurVioletPale}{violet!10!white}%
\newcommand{\seqCouleurVertModere}{green!50!white}%
\newcommand{\seqCouleurVertPale}{green!10!white}%
%% ---- Algorithmique et programmation
\newcommand{\seqCouleurRoseVif}{magenta!85!black}%
\newcommand{\seqCouleurRoseMoyen}{magenta!25!white}%
\newcommand{\seqCouleurRoseModere}{magenta!50!white}%
\newcommand{\seqCouleurRosePale}{magenta!10!white}%
\newcommand{\seqCouleurLimeModere}{lime!50!white}%
\newcommand{\seqCouleurLimePale}{lime!10!white}%
%% ---- Proportionnalité (cycle 3, thème E) -- v0.52 : canard vif, orange pâle
\newcommand{\seqCouleurCanardVif}{teal!85!black}%
\newcommand{\seqCouleurCanardMoyen}{teal!25!white}%
\newcommand{\seqCouleurCanardModereVif}{teal!50!white}%
\newcommand{\seqCouleurCanardTresPale}{teal!10!white}%
%% ---- Autre
\newcommand{\seqCouleurNoirVif}{black!85!black}%
\newcommand{\seqCouleurNoirMoyen}{black!25!white}%
\newcommand{\seqCouleurNoirModere}{black!50!white}%
\newcommand{\seqCouleurNoirPale}{black!10!white}%
\newcommand{\seqCouleurGrisModere}{gray!50!white}%
\newcommand{\seqCouleurGrisPale}{gray!10!white}%

% -------------------------------------------------------------------------------------------
% Règle les deux couleurs avec leurs niveaux d'intensité pour le thème spécifié
% #1 : nom du code couleurs
\newcommand{\seq@setColors}[1]{%
	\ifthenelse{ \equal{#1}{nombres} }%
	{%
		\def\seq@couleuri@vif{\seqCouleurRougeVif}%
		\def\seq@couleuri@moyen{\seqCouleurRougeMoyen}%
		\def\seq@couleuri@modere{\seqCouleurRougeModere}%
		\def\seq@couleuri@pale{\seqCouleurRougePale}%
		\def\seq@couleurii@modere{\seqCouleurJauneModere}%
		\def\seq@couleurii@pale{\seqCouleurJaunePale}%
	}%
	{%
		\ifthenelse{ \equal{#1}{donnees} }%
		{%
			\def\seq@couleuri@vif{\seqCouleurBleuVif}%
			\def\seq@couleuri@moyen{\seqCouleurBleuMoyen}%
			\def\seq@couleuri@modere{\seqCouleurBleuModere}%
			\def\seq@couleuri@pale{\seqCouleurBleuPale}%
			\def\seq@couleurii@modere{\seqCouleurCanardModere}%
			\def\seq@couleurii@pale{\seqCouleurCanardPale}%
		}%
		{%
			\ifthenelse{ \equal{#1}{grandeurs} }%
			{%
				\def\seq@couleuri@vif{\seqCouleurMarronVif}%
				\def\seq@couleuri@moyen{\seqCouleurMarronMoyen}%
				\def\seq@couleuri@modere{\seqCouleurMarronModere}%
				\def\seq@couleuri@pale{\seqCouleurMarronPale}%
				\def\seq@couleurii@modere{\seqCouleurOrangeModere}%
				\def\seq@couleurii@pale{\seqCouleurOrangePale}%
			}%
			{%
				\ifthenelse{ \equal{#1}{geometrie} }%
				{%
					\def\seq@couleuri@vif{\seqCouleurVioletVif}%
					\def\seq@couleuri@moyen{\seqCouleurVioletMoyen}%
					\def\seq@couleuri@modere{\seqCouleurVioletModere}%
					\def\seq@couleuri@pale{\seqCouleurVioletPale}%
					\def\seq@couleurii@modere{\seqCouleurVertModere}%
					\def\seq@couleurii@pale{\seqCouleurVertPale}%
				}%
				{%
					\ifthenelse{ \equal{#1}{algorithmique} }%
					{%
						\def\seq@couleuri@vif{\seqCouleurRoseVif}%
						\def\seq@couleuri@moyen{\seqCouleurRoseMoyen}%
						\def\seq@couleuri@modere{\seqCouleurRoseModere}%
						\def\seq@couleuri@pale{\seqCouleurRosePale}%
						\def\seq@couleurii@modere{\seqCouleurLimeModere}%
						\def\seq@couleurii@pale{\seqCouleurLimePale}%
					}%
					{%
					\ifthenelse{ \equal{#1}{proportionnalite} }%
					{%
						\def\seq@couleuri@vif{\seqCouleurCanardVif}%
						\def\seq@couleuri@moyen{\seqCouleurCanardMoyen}%
						\def\seq@couleuri@modere{\seqCouleurCanardModereVif}%
						\def\seq@couleuri@pale{\seqCouleurCanardTresPale}%
						\def\seq@couleurii@modere{\seqCouleurOrangeModere}%
						\def\seq@couleurii@pale{\seqCouleurOrangePale}%
					}%
					{%
						\def\seq@couleuri@vif{\seqCouleurNoirVif}%
						\def\seq@couleuri@moyen{\seqCouleurNoirMoyen}%
						\def\seq@couleuri@modere{\seqCouleurNoirModere}%
						\def\seq@couleuri@pale{\seqCouleurNoirPale}%
						\def\seq@couleurii@modere{\seqCouleurGrisModere}%
						\def\seq@couleurii@pale{\seqCouleurGrisPale}%
					}%
					}%
				}%
			}%
		}%
	}%
}%
% -------------------------------------------------------------------------------------------

%% ---- Couleur des lignes de séparation dans les tableaux
\newcommand{\seq@rule@color}[1]{ \def\CT@arc@{\color{#1}} }%
% -------------------------------------------------------------

%% ---- Couleur du tableau d'objectifs
\newcolumntype{Y}{>{\columncolor{\seq@couleurii@pale}}X}%
%    Types de colonnes
\newcolumntype{L}[1]{>{\raggedright\arraybackslash}p{#1}}%
\newcolumntype{C}[1]{>{\centering\arraybackslash}p{#1}}%

%% ------------------------------------------------------------------
%% ---- Macro pour fiches à compléter -------------------------------
%% ------------------------------------------------------------------
% \begin{macro}{\acompleter}%
% \IndexMacro{\acompleter}%
% Crée un espace souligné pour compléter (utilisé dans les flashcards)
% #1 : texte fantôme (détermine la longueur de l'espace)
% \begin{macrocode}%
\newcommand{\acompleter}[1]{\;\underline{\phantom{#1}}\;}%
% \end{macrocode}%
% \end{macro}%

%% ------------------------------------------------------------------
%% ---- Encadrés colorés pour mise en valeur -----------------------
%% ------------------------------------------------------------------
% \begin{macro}{\seqBoiteJaune}%
% \IndexMacro{\seqBoiteJaune}%
% Encadré avec fond coloré pour mise en valeur dans le cours
% Utilise la couleur secondaire modérée du thème courant
% #1 : contenu à encadrer
% \begin{macrocode}%
\newcommand{\seqBoiteJaune}[1]{%
    \adjustbox{cframe={\seq@couleurii@modere} 2pt 2pt 2pt,bgcolor=\seq@couleurii@modere}{#1}%
}%
% \end{macrocode}%
% \end{macro}%
%% Alias historique — certains livrets appellent \boiteJauneModere
%% (même effet, nom non préfixé seq). Maintenu pour compatibilité.
\newcommand{\boiteJauneModere}[1]{%
    \adjustbox{cframe={\seqCouleurJauneModere} 2pt 2pt 2pt,bgcolor=\seqCouleurJauneModere}{#1}%
}%

% -------------------------------------------------------------
% Styles d'encadrés et d'items visuels récurrents
% (ex : boîtes tcolorbox pour Objectifs / Prérequis / Savoir-faire)
% -------------------------------------------------------------
%
% \begin{environment}{seqBlocObjectifs}%
% \IndexEnv{seqBlocObjectifs}%
% Encart visuel stylé pour présenter les objectifs.
%
% \begin{macrocode}%
% \newenvironment{seqBlocObjectifs}{%
%   % tcolorbox / minipage / tikz node / etc.
% }{%
%   % fin
% }%
% \end{macrocode}%
% \end{environment}%
%
% Idem pour blocs "Pré-requis", "Savoir-faire", "Méthodes", etc.
%
%% ------------------------------------------------------------------
%% ---- Différentes boites ------------------------------------------
%% ------------------------------------------------------------------
\newtcolorbox{boiteTitreGen}[1]{%
	title=#1,
	colback=\seq@couleuri@pale,
	colframe=\seq@couleuri@vif,
	coltext=\seq@couleuri@vif,
	colbacktitle=\seq@couleuri@vif,
	enhanced,
	attach boxed title to top left={yshift=-5pt,xshift=20pt}%
}%
\newtcolorbox{boiteTitreSection}{%
	colback=\seq@couleuri@pale,
	colframe=\seq@couleuri@vif,
	coltext=\seq@couleuri@vif
}%
\newtcolorbox{boiteTitreSousSection}{%
	colback=\seq@couleuri@pale,
	colframe=\seq@couleuri@vif,
	coltext=\seq@couleuri@vif,
	width=0.6\linewidth
}%
\newtcolorbox{boiteTitreSerie}{%
	colback=\seq@couleurii@modere,
	colframe=\seq@couleuri@vif,
	coltext=\seq@couleuri@vif
}%
\newtcolorbox{boiteTitreSerieCorrige}{%
	colback=\seq@couleuri@pale,
	colframe=\seq@couleuri@vif,
	coltext=\seq@couleuri@vif,
	width=0.5\linewidth
}%
\newtcolorbox{boiteTitreGroupe}{%
	colback=\seq@couleurii@modere,
	colframe=\seq@couleuri@vif,
	coltext=\seq@couleuri@vif
}%
\newtcolorbox{boiteTitreExo}{%
	colback=\seq@couleuri@pale,
	colframe=\seq@couleuri@vif,
	coltext=\seq@couleuri@vif
}%
\newtcolorbox{boiteNumExo}{%
	width=100pt,
	colback=\seq@couleuri@pale,
	colframe=\seq@couleuri@vif,
	coltext=\seq@couleuri@vif,
	after={},
	box align=center,
	boxsep=2pt,
	halign=right
}%
\newtcolorbox{boitePaleNoBreak}[1]{%
	title=#1,
	colback=\seq@couleurii@pale,
	colframe=\seq@couleuri@vif,
	colbacktitle=\seq@couleuri@vif,
	fonttitle=\bfseries,
	enhanced,
	attach boxed title to top left={yshift=-5pt,xshift=20pt}%
}%
\newtcolorbox{boitePaleBreak}[1]{%
	title=#1,
	colback=\seq@couleurii@pale,
	colframe=\seq@couleuri@vif,
	colbacktitle=\seq@couleuri@vif,
	breakable,
	fonttitle=\bfseries,
	enhanced,
	attach boxed title to top left={yshift=-5pt,xshift=20pt},
	height fixed for=all,
	pad after break=20pt
}%
\newtcolorbox{boitePaleBreakFitHeight}[1]{%
	title=#1,
	colback=\seq@couleurii@pale,
	colframe=\seq@couleuri@vif,
	colbacktitle=\seq@couleuri@vif,
	breakable,
	fonttitle=\bfseries,
	enhanced,
	attach boxed title to top left={yshift=-5pt,xshift=20pt},
	pad after break=20pt
}%
\newtcolorbox{boitePaleBreakNoTitle}{%
	colback=\seq@couleurii@pale,
	colframe=\seq@couleuri@vif,
	breakable,
	enhanced,
	height fixed for=all
}%
\newtcolorbox{boitePaleNoBreakNoTitle}{%
	colback=\seq@couleurii@pale,
	colframe=\seq@couleuri@vif
}%
\newtcolorbox{boitePaleVariable}[1]{%
	colback=\seq@couleurii@pale,
	colframe=\seq@couleuri@vif,
	width=#1
}%
\newtcolorbox{boitePaleVariableNobeforeafter}[1]{%
	colback=\seq@couleurii@pale,
	colframe=\seq@couleuri@vif,
	nobeforeafter,
	width=#1
}%
\newtcolorbox{boiteVariableTitre}[2]{%
	title={\footnotesize #1},
	colback=white,
	colframe=\seq@couleuri@vif,
	colbacktitle=\seq@couleuri@vif,
	fonttitle=\bfseries,
	enhanced,
	attach boxed title to top left={yshift=-5pt,xshift=20pt},
	width=#2
}%
\newtcolorbox{boiteVariableWidthHeightTitre}[3]{%
	title={\footnotesize #1},
	colback=\seq@couleurii@pale,
	colframe=\seq@couleuri@vif,
	colbacktitle=\seq@couleuri@vif,
	fonttitle=\bfseries,
	enhanced,
	attach boxed title to top left={yshift=-5pt,xshift=20pt},
	width=#2, height=#3
}%
\newtcolorbox{boitePaleNoBreakHeightNoTitle}[1]{%
	colback=\seq@couleurii@pale,
	colframe=\seq@couleuri@vif,
	height=#1
}%

% \begin{macro}{seqRenduTableauPrerequis}%
% \IndexMacro{\seqRenduTableauPrerequis}%
% Génère la boîte « Prérequis » avec son contenu, pour le livret de séquence
% (cf. \texttt{services/livret\_sequence.py:\_generer\_tableau\_prerequis}).
%
% Contrairement à \texttt{seqRenduTableauObjectifs}, le contenu n'est pas un
% tableau mais un bloc de texte libre (typiquement des
% \texttt{\textbackslash subsection*} suivis d'\texttt{itemize} listant
% les objectifs des séquences précédentes).
%
% La boîte hérite des couleurs du thème courant
% (\texttt{\textbackslash seqSetColorsTheme}). Saut de page autorisé
% (\texttt{breakable}).
%
% Usage : \texttt{\textbackslash seqRenduTableauPrerequis\{contenu\}}%
%
%    \begin{macrocode}%
\newtcolorbox{seq@boite@prerequis}{%
	title=Prérequis,
	colback=\seq@couleurii@pale,
	colframe=\seq@couleuri@vif,
	colbacktitle=\seq@couleuri@vif,
	breakable,
	fonttitle=\bfseries,
	enhanced,
	attach boxed title to top left={yshift=-5pt,xshift=20pt},
	pad after break=20pt
}%
\newcommand{\seqRenduTableauPrerequis}[1]{%
	\begin{seq@boite@prerequis}%
		#1%
	\end{seq@boite@prerequis}%
}%
%    \end{macrocode}%
% \end{macro}%

% \begin{macro}{seqRenduTableauObjectifs}%
% \IndexMacro{\seqRenduTableauObjectifs}%
% Génère la boîte « Objectifs » avec son tableau intégré, pour le livret
% de séquence (cf.
% \texttt{services/livret\_sequence.py:\_generer\_tableau\_objectifs}).
%
% Le tableau est intégré dans la boîte via la clé tcolorbox
% \texttt{tabularx*={}\{X|X|X|X\}} : pas de
% \texttt{\textbackslash begin\{tabularx\}} séparé à émettre côté code.
% Le contenu \texttt{\#1} est constitué uniquement des lignes de données
% (chaque ligne \texttt{Objectif \& F \& A \& E \textbackslash\textbackslash\textbackslash hline}).
%
% L'en-tête fixe est inclus dans la définition même de la commande, en
% deux niveaux :
% \begin{itemize}%
% \item ligne 1 : cellule vide (sur 2 lignes via \texttt{\textbackslash multirow}) puis
%   « Niveau de maîtrise » couvrant les 3 colonnes critères
%   (\texttt{\textbackslash multicolumn})
% \item ligne 2 : « À consolider », « Satisfaisant », « Très bien »
% \end{itemize}%
%
% La boîte hérite des couleurs du thème courant
% (\texttt{\textbackslash seqSetColorsTheme}). Saut de page autorisé.
%
% Le type de colonne \texttt{C} (raggedright avec \texttt{arraybackslash})
% est défini localement pour permettre au \texttt{\textbackslash multicolumn}%
% de gérer correctement la largeur via \texttt{\textbackslash hsize}.
%
% Usage : \texttt{\textbackslash seqRenduTableauObjectifs\{ligne1 \textbackslash\textbackslash\textbackslash hline ligne2 \textbackslash\textbackslash\textbackslash hline …\}}%
%
%    \begin{macrocode}%
\newtcolorbox{seq@boite@objectifs}[1]{%
	title={#1},
	colback=\seq@couleurii@pale,
	colframe=\seq@couleuri@vif,
	colbacktitle=\seq@couleuri@vif,
	breakable,
	fonttitle=\bfseries,
	enhanced,
	attach boxed title to top left={yshift=-5pt,xshift=20pt},
	pad after break=20pt,
	tabularx*={}{X|X|X|X}%
}%
\newcommand{\seqRenduTableauObjectifs}[2][Objectifs]{%
	\begingroup
	\newcolumntype{C}{>{\centering\arraybackslash}X}%
	\begin{seq@boite@objectifs}{#1}%
		\multirow{2}{\hsize}{} & \multicolumn{3}{|>{\hsize=3\hsize}C|}{Niveau de maîtrise} \\\cline{2-4}%
		& \`A consolider & Satisfaisant & Très bien \\\hline
		#2
	\end{seq@boite@objectifs}%
	\endgroup
}%
%    \end{macrocode}%
% \end{macro}%

%% ------------------------------------------------------------------
%% ---- Cellules dans les tableaux d'objectif -----------------------
%% ------------------------------------------------------------------
\newcommand\cellVertUn[1]{%
	\multirow{4}{*}{%
		\begin{minipage}{1em}%
			\adjustbox{angle=90,valign=M}{\color{\seq@couleuri@vif} \textbf{#1}}%
		\end{minipage}%
	}%
}%
\newcommand\cellVert[1]{%
	\adjustbox{angle=90,valign=M}{\color{\seq@couleuri@vif} \textbf{#1}}%
}%
\newcommand{\seq@headcol}{%
	\rowcolor{\seq@couleurii@pale}%
}%
\newcommand{\seq@headtextcol}{%
	\color{\seq@couleuri@vif}%
}%
% -------------------------------------------------------------

%% ------------------------------------------------------------------
%% ---- Pour les flashcards -----------------------------------------
%% ------------------------------------------------------------------
\newcommand{\seqBoiteFlashcardVide}{%
	\begin{tcolorbox}[title=\begin{minipage}{0.8\textwidth}\Large\phantom{Vide!}\end{minipage}%
			,height fill
			,colback=\seq@couleurii@pale
			,colframe=\seq@couleuri@vif
			,colbacktitle=\seq@couleuri@vif
			,fonttitle=\bfseries
			,enhanced
			,attach boxed title to top left={yshift=-5pt,xshift=20pt}%
		]
	\end{tcolorbox}%
	\newpage
}%
% -------------------------------------------------------------------------------------------
% \begin{macro}{\seqBoiteTitreFlashcard}%
% \IndexMacro{\seqBoiteTitreFlashcard}%
% Met en forme le titre de la flashcard de l'objectif spécifié
% #1 : code objectif
% #2 : complément au titre de la flashcard (Ex: 1\iere{} fiche}%
% \begin{macrocode}%
\newcommand{\seqBoiteTitreFlashcard}[2]{%
	\stepcounter{flashcardnb}%
	\begin{tcolorbox}[title=\begin{minipage}{0.8\textwidth}\Large \seq@code@sequence: \seqSequenceGetNom\end{minipage}%
			,height fill
			,colback=\seq@couleurii@pale
			,colframe=\seq@couleuri@vif
			,colbacktitle=\seq@couleuri@vif
			,fonttitle=\bfseries
			,enhanced
			,attach boxed title to top left={yshift=-5pt,xshift=20pt}%
		]
		\begin{minipage}[t][9em]{\textwidth-20pt}Classe de \seqNiveauGetNomCourt{\seq@code@niveau}\hfill Fiche numéro \theflashcardnb\end{minipage}\par
		\begin{minipage}[c][7em]{\textwidth-20pt}\begin{center}\huge\seqObjectifGetNom{#1}\end{center}\end{minipage}\par
		\begin{minipage}[b][6em]{\textwidth-20pt}\hfill #2\end{minipage}\par
		\begin{minipage}[b][6em]{\textwidth-20pt}Version de \datemoisannee\today \hfill Nom et prénom: \acompleter{un très long nom et un prénom aussi}\end{minipage}%
	\end{tcolorbox}%
	\newpage
}%
% \end{macrocode}%
% \end{macro}%
% -------------------------------------------------------------------------------------------
%% Options
%% Environnement contenu de flashcard
\define@cmdkey[seq]{boiteContenuFlashcard}{titre}[Exemples]{}%
\newenvironment{seqBoiteContenuFlashcard}[1][]{%
	\setkeys[seq]{boiteContenuFlashcard}{titre,#1}%
	\begin{tcolorbox}[title=\cmdseq@boiteContenuFlashcard@titre
			,natural height
			,colback=\seq@couleurii@pale
			,colframe=\seq@couleuri@vif
			,colbacktitle=\seq@couleuri@vif
			,fonttitle=\bfseries
			,enhanced
			,attach boxed title to top left={yshift=-5pt,xshift=20pt}%
		]
		}{%
	\end{tcolorbox}%
}%
%% Options
\define@cmdkey[seq]{boiteFillContenuFlashcard}{titre}[Exemples]{}%
%% Environnement contenu de flashcard
\newenvironment{seqBoiteFillContenuFlashcard}[1][]{%
	\setkeys[seq]{boiteFillContenuFlashcard}{titre,#1}%
	\begin{tcolorbox}[title=\cmdseq@boiteFillContenuFlashcard@titre
			,height fill
			,colback=\seq@couleurii@pale
			,colframe=\seq@couleuri@vif
			,colbacktitle=\seq@couleuri@vif
			,fonttitle=\bfseries
			,enhanced
			,attach boxed title to top left={yshift=-5pt,xshift=20pt}%
		]
		}{%
	\end{tcolorbox}%
	\newpage
}%

%% ------------------------------------------------------------------
%% ---- Blocs de niveau 1 et 2 --------------------------------------
%% ------------------------------------------------------------------
%    \begin{environment}{seqBloc}%
\newenvironment{seqBloc}[1]{%
	\begin{boiteTitreSection}%
		\begin{minipage}[c][12pt][c]{\linewidth}%
			\centering \large \textbf{#1}%
		\end{minipage}%
	\end{boiteTitreSection}%
	\begin{boitePaleNoBreakNoTitle}%
		\begin{minipage}{\linewidth}%
			\begin{dinglist}{224}%
				}{%
			\end{dinglist}%
		\end{minipage}%
	\end{boitePaleNoBreakNoTitle}%
}%
%    \end{environment}%

%    \begin{environment}{seqBlocN2}%
\newenvironment{seqBlocN2}{%
	\begin{dinglist}{118}%
		}{%
	\end{dinglist}%
}%
%    \end{environment}%

% -----------------------------------------------
% Environnement enumerate sur plusieurs colonnes
% -----------------------------------------------
%    \begin{environment}{colEnum}%
%% Options
\define@cmdkey[seq]{colEnum}{nbCols}[2]{}%
\define@cmdkey[seq]{colEnum}{label}[\arabic*)]{}%
\newenvironment{seqColEnum}[1][]{%
	\setkeys[seq]{colEnum}{nbCols,label,#1}%
	\removelastskip
	\ifthenelse{\cmdseq@colEnum@nbCols > 1}%
	{\begin{multicols}{\cmdseq@colEnum@nbCols}}{}%
			\begin{enumerate}[label=\cmdseq@colEnum@label]
				}{%
			\end{enumerate}%
			\ifthenelse{\cmdseq@colEnum@nbCols > 1}{\end{multicols}}{}%
}%
%    \end{environment}%

% -----------------------------------------------
% Environnement itemize sur plusieurs colonnes
% -----------------------------------------------
%    \begin{environment}{colItem}%
%% Options
\define@cmdkey[seq]{colItem}{nbCols}[2]{}%
\define@cmdkey[seq]{colItem}{label}[$\bullet$]{}%
\newenvironment{seqColItem}[1][]{%
	\setkeys[seq]{colItem}{nbCols,label,#1}%
	\ifthenelse{\cmdseq@colItem@nbCols > 1}%
	{\begin{multicols}{\cmdseq@colItem@nbCols}}{}%
			\begin{itemize}[label=\cmdseq@colItem@label]
				}{%
			\end{itemize}%
			\ifthenelse{\cmdseq@colItem@nbCols > 1}{\end{multicols}}{}%
}%
%    \end{environment}%

%% ------------------------------------------------------------------
%% ---- Pour Tikz ---------------------------------------------------
%% ------------------------------------------------------------------
\tikzstyle{axe}=[line width=0.3mm, ->, color=black, line cap=rect]

%    \begin{macro}{\seqOrigine}%
%    \changes{v1.0}{2018/05/02}{Version initiale}%
\newcommand{\seqOrigine}{\draw (0,0) node[below left] {\small $0$};}%
%    \end{macro}%

%    \begin{macro}{\seqAxeX}%
%    \changes{v1.0}{2018/05/02}{Version initiale}%
%% Options
\define@cmdkey[seq]{axeX}{ord}[0]{}%
\define@cmdkey[seq]{axeX}{absMin}{}%
\define@cmdkey[seq]{axeX}{absMax}{}%
\define@cmdkey[seq]{axeX}{absList}{}%
\define@cmdkey[seq]{axeX}{libelle}{}%
\newcommand{\seqAxeX}[1][]{%
	\setkeys[seq]{axeX}{ord,#1}%
	\draw[axe] (\cmdseq@axeX@absMin,\cmdseq@axeX@ord)--
	(\cmdseq@axeX@absMax,\cmdseq@axeX@ord)
	\ifx\cmdseq@axeX@libelle\undefined\else
		node[above] {\textit{\cmdseq@axeX@libelle}}%
	\fi;
	\foreach \x in \cmdseq@axeX@absList {%
		\draw[thin] (\x,\cmdseq@axeX@ord) node {\small $+$};
		\draw (\x,\cmdseq@axeX@ord) node[below] {\small $\x$};
	}%
}%
%    \end{macro}%

%    \begin{macro}{\seqAxeY}%
%    \changes{v1.0}{2018/05/02}{Version initiale}%
%% Options
\define@cmdkey[seq]{axeY}{abs}[0]{}%
\define@cmdkey[seq]{axeY}{ordMin}{}%
\define@cmdkey[seq]{axeY}{ordMax}{}%
\define@cmdkey[seq]{axeY}{ordList}{}%
\define@cmdkey[seq]{axeY}{libelle}{}%
\newcommand{\seqAxeY}[1][]{%
	\setkeys[seq]{axeY}{abs,#1}%
	\draw[axe] (\cmdseq@axeY@abs,\cmdseq@axeY@ordMin)--
	(\cmdseq@axeY@abs,\cmdseq@axeY@ordMax)
	\ifx\cmdseq@axeY@libelle\undefined\else
		node[right] {\textit{\cmdseq@axeY@libelle}}%
	\fi;
	\foreach \y in \cmdseq@axeY@ordList {%
		\draw[thin] (\cmdseq@axeY@abs, \y) node {\small $+$};
		\draw (\cmdseq@axeY@abs, \y) node[left] {\small $\y$};
	}%
}%
%    \end{macro}%

%% ------------------------------------------------------------------
%% ---- Réglages généraux -------------------------------------------
%% ------------------------------------------------------------------
%% ---- Format de date
\newdateformat{datemoisannee}{\monthname\;\THEYEAR}%
%% ---- Espacement autour des fractions
\newcommand\TopStrut{\rule{0pt}{2.6ex}}       % Top strut
\newcommand\BottomStrut{\rule[-1.2ex]{0pt}{0pt}} % Bottom strut
%% ---- Expansion des options dans les fichiers de corrigé
\newcommand\eue{\expandafter\unexpanded\expandafter}%
%% ---- Taille des marges
\ifdocstd \geometry{left=35pt,right=35pt,top=50pt,bottom=70pt,head=30pt,foot=10pt}%
\else \geometry{left=100pt,right=50pt,top=70pt,bottom=70pt,head=30pt,foot=10pt}%
\fi
%% ---- Pas d'indentation des paragraphes
\setlength{\parindent}{0pt}%
%% ---- Espacement avant le pied de page
\setlength{\footskip}{15.0pt}%
%% ---- Pas de trait sous l'entête
\renewcommand{\headrulewidth}{0pt}%
%% ---- Pas de trait au dessus du pied de page
\renewcommand{\footrulewidth}{0pt}%


%% ------------------------------------------------------------------
%% ---- Pour tabularray ---------------------------------------------
%% ------------------------------------------------------------------
%% ---- Texte de continuation des longs tableaux
\DefTblrTemplate{contfoot-text}{normal}{Suite page suivante}%
\SetTblrTemplate{contfoot-text}{normal}%
\DefTblrTemplate{conthead-text}{normal}{(suite)}%
\SetTblrTemplate{conthead-text}{normal}%

%% ---- Titre vertical de colonne (tableaux étroits à entête long)
%% Usage : \seqTitreTabV{Libellé long} dans une cellule d'entête.
\newcommand{\seqTitreTabV}[1]{%
    \rotatebox{70}{\begin{minipage}{8em}#1\end{minipage}}%
}%

%% ------------------------------------------------------------------
%% ---- Pour multicols ----------------------------------------------
%% ------------------------------------------------------------------
%% ---- Pas de ligne verticale entre les colonnes
\renewcommand{\columnseprule}{0pt}%
%% ---- Séparation par défaut entre les colonnes
\setlength{\columnsep}{10pt}%
%% ---- Espace avant et après multicols réduit de moitié
\setlength{\multicolsep}{6.0pt plus 2.0pt minus 1.5pt}%


% -------------------------------------------------------------
% Icônes / pictos
% Si tu as des macros genre \seqPictoObjectif, \seqPictoInfo...
% c'est ici.
% -------------------------------------------------------------
% \begin{macro}{\seqPictoObjectif}%
% \IndexMacro{\seqPictoObjectif}%
% \begin{macrocode}%
% \newcommand{\seqPictoObjectif}{...}%
% \end{macrocode}%
% \end{macro}%


% =============================================================================
% Couleurs pour les fiches PPRE (Plan Personnalisé de Réussite Éducative)
% =============================================================================
%
% Les couleurs des trois séries F / A / E sont indépendantes du thème
% mathématique courant (domaine). Elles encodent un niveau de réussite,
% pas un domaine.
%
% Correspondance avec le système de notation du paquet :
%   F (fondamentale) = à consolider  → ambre/orange
%   A (avancée)      = satisfaisant  → bleu clair
%   E (exploration)  = très bon      → vert clair
%
% \begin{macrocode}%
%% ---- Série Fondamentale (F) : ambre — à consolider
\newcommand{\seqCouleurSerieF}{orange!18!white}%
\newcommand{\seqCouleurSerieFCadre}{orange!65!black}%
%% ---- Série Avancée (A) : bleu clair — satisfaisant
\newcommand{\seqCouleurSerieA}{blue!12!white}%
\newcommand{\seqCouleurSerieACadre}{blue!55!black}%
%% ---- Série Exploration (E) : vert clair — très bon
\newcommand{\seqCouleurSerieE}{green!15!white}%
\newcommand{\seqCouleurSerieECadre}{green!50!black}%
% \end{macrocode}%
%
% Boîte interne pour les sections F/A/E d'une fiche PPRE.
% \begin{macro}{seqPPRE@boite}%
% \begin{macrocode}%
\newtcolorbox{seqPPRE@boite}[3]{%
    title          = {#1},%
    colback        = {#2},%
    colframe       = {#3},%
    colbacktitle   = {#3},%
    coltitle       = white,%
    fonttitle      = \bfseries\small,%
    enhanced,%
    attach boxed title to top left={yshift=-4pt, xshift=8pt},%
    top            = 8pt,%
    bottom         = 4pt,%
    left           = 4pt,%
    right          = 4pt,%
    before skip    = 6pt,%
    after skip     = 4pt,%
}%
% \end{macrocode}%
% \end{macro}%

% \begin{macro}{\seqStyleObjectifsHeader}%
% Wrapper public pour styliser un en-tête de cellule en couleur du thème
% courant. Mise au point en v0.11.6.2 pour permettre au tableau d'objectifs
% inliné côté appli (cf. livret_sequence.py:_generer_tableau_objectifs)
% d'utiliser la couleur du thème sans passer par la macro \seq@couleuri@vif
% qui contient un caractère @ et n'est donc pas utilisable depuis un .tex
% sans \makeatletter.
%
% Usage : \seqStyleObjectifsHeader{Texte de l'en-tête}%
%
% Le texte est rendu en gras et en couleur \seq@couleuri@vif (la couleur
% « vif » du thème courant, fixée par \seqSetColorsTheme).
%    \begin{macrocode}%
\newcommand{\seqStyleObjectifsHeader}[1]{%
    \color{\seq@couleuri@vif}\textbf{#1}%
}%
%    \end{macrocode}%
% \end{macro}%

% Fin du module theme
%</package>
